\section*{الفصل \ref{ch:g9:sqrt} --- الجذور التربيعية}

\begin{solution}{exo:g9:sqrt:1}
$\sqrt{64} = 8$; \quad $\sqrt{100} = 10$; \quad
$\sqrt{0.09} = 0.3$ (لأنّ $0.3^2 = 0.09$)؛ \quad
$\sqrt{\frac{1}{25}} = \frac15$; \quad
$\left(\sqrt{13}\right)^2 = 13$; \quad
$\sqrt{(-4)^2} = \sqrt{16} = 4 = \abs{-4}$.
\end{solution}

\begin{solution}{exo:g9:sqrt:2}
$\sqrt{50} = \sqrt{25 \times 2} = 5\sqrt2$; \quad
$\sqrt{45} = \sqrt{9 \times 5} = 3\sqrt5$; \quad
$\sqrt{98} = \sqrt{49 \times 2} = 7\sqrt2$; \quad
$\sqrt{300} = \sqrt{100 \times 3} = 10\sqrt3$.
\end{solution}

\begin{solution}{exo:g9:sqrt:3}
$\sqrt2 \times \sqrt{18} = \sqrt{36} = 6$.

$\dfrac{\sqrt{75}}{\sqrt3} = \sqrt{\dfrac{75}{3}} = \sqrt{25} = 5$.

$\sqrt5 \times \sqrt{20} = \sqrt{100} = 10$.

$\dfrac{\sqrt8}{\sqrt{50}} = \sqrt{\dfrac{8}{50}} = \sqrt{\dfrac{4}{25}}
= \dfrac25$.
\end{solution}

\begin{solution}{exo:g9:sqrt:4}
$x^2 = 36$: $x = 6$ أو $x = -6$.

$x^2 = 11$: $x = \sqrt{11}$ أو $x = -\sqrt{11}$.

$x^2 + 4 = 0$ يعني $x^2 = -4 < 0$: لا حلّ.

$2x^2 = 50$: $x^2 = 25$، إذن $x = 5$ أو $x = -5$.
\end{solution}

\begin{solution}{exo:g9:sqrt:5}
$\sqrt{20} + \sqrt{45} = 2\sqrt5 + 3\sqrt5 = 5\sqrt5$.

$3\sqrt8 - \sqrt{18} + \sqrt2 = 3 \times 2\sqrt2 - 3\sqrt2 + \sqrt2
= (6 - 3 + 1)\sqrt2 = 4\sqrt2$.
\end{solution}

\begin{solution}{exo:g9:sqrt:6}
$\left(\sqrt3 + 1\right)^2 = 3 + 2\sqrt3 + 1 = 4 + 2\sqrt3$.

$\left(2\sqrt5 - 3\right)^2 = 4 \times 5 - 2 \times 3 \times 2\sqrt5 + 9
= 29 - 12\sqrt5$.

$\left(\sqrt6 + \sqrt2\right)\left(\sqrt6 - \sqrt2\right) = 6 - 2 = 4$.
\end{solution}

\begin{solution}{exo:g9:sqrt:7}
الضلع: $\sqrt{6400} = 80$ م. \omterm{def:g4:geometry:circle}{والقطر} (فيثاغورس):
$\sqrt{80^2 + 80^2} = 80\sqrt2 \approx 113$ م.
\end{solution}

\begin{solution}{exo:g9:sqrt:8}
$\dfrac{1}{\sqrt2} = \dfrac{1 \times \sqrt2}{\sqrt2 \times \sqrt2}
= \dfrac{\sqrt2}{2}$.

$\dfrac{6}{\sqrt3} = \dfrac{6\sqrt3}{3} = 2\sqrt3$.
\quad
$\dfrac{10}{\sqrt5} = \dfrac{10\sqrt5}{5} = 2\sqrt5$.
\end{solution}

\begin{solution}{exo:g9:sqrt:9}
\emph{1. صواب}: وهي قاعدة \omterm{def:g2:mult:def}{الجداء} (\cref{thm:g9:sqrt:rules}).

\emph{2. خطأ}: $\sqrt{9 + 16} = 5$ لكن $\sqrt9 + \sqrt{16} = 7$.

\emph{3. خطأ} من أجل $x$ سالب: $\sqrt{(-4)^2} = 4 \neq -4$.
والصياغة الصحيحة هي $\sqrt{x^2} = \abs{x}$.

\emph{4. صواب}: $\left(3\sqrt2\right)^2 = 9 \times 2 = 18$.
\end{solution}

\begin{solution}{exo:g9:sqrt:10}
$\left(2 + \sqrt3\right)^2 = 4 + 4\sqrt3 + 3 = 7 + 4\sqrt3$. إذن
$x = \sqrt{7 + 4\sqrt3} = \sqrt{\left(2+\sqrt3\right)^2} = 2 + \sqrt3$
(وهو عدد غير سالب، فيزيل الجذر المربّع فحسب).

وبالمثل $\left(\sqrt5 - 2\right)^2 = 5 - 4\sqrt5 + 4 = 9 - 4\sqrt5$،
و $\sqrt5 - 2 > 0$، إذن
$\sqrt{9 - 4\sqrt5} = \sqrt5 - 2$ (لا $2 - \sqrt5$، فهو
سالب!).
\end{solution}

\begin{solution}{pb:g9:sqrt:1}
\textbf{1.} $1.4^2 = 1.96 < 2 < 2.25 = 1.5^2$؛
و $1.41^2 = 1.9881 < 2 < 2.0164 = 1.42^2$. وبثلاثة أرقام عشرية:
$1.414^2 = 1.999396 < 2 < 2.002225 = 1.415^2$، إذن
$1.414 < \sqrt2 < 1.415$.

\textbf{2.} بتربيع $\sqrt2 = \frac ab$ نجد
$2 = \frac{a^2}{b^2}$، ومن ثمّ $a^2 = 2b^2$: فالعدد $a^2$ ضعف عدد
صحيح --- أي \emph{زوجي}.

\textbf{3.} لو كان $a$ فرديًا، لكان $a^2$ فرديًا
(\cref{pb:g8:pythagoras:1})؛ وبما أنّ $a^2$ زوجي، فالعدد $a$
زوجي: أي $a = 2k$. وبالتعويض: $4k^2 = 2b^2$، إذن $b^2 = 2k^2$
زوجي، وبواقعة الزوجية نفسها يكون $b$ زوجيًا.

\textbf{4.} كون $a$ و $b$ زوجيين معناه أنّهما \omterm{def:g6:wholes:divisible}{قابلان للقسمة}
على $2$ --- لكنّ $\frac ab$ كان مختزلًا إلى أبسط صورة، فلا يشتركان
في أيّ قاسم. وهذا تناقض: فلا يمكن أن يوجد \omterm{def:g9:fractions:fraction}{الكسر} المفترض.
\emph{المبرهنة: العدد $\sqrt2$ أصمّ.}

\textbf{5.} يقيم التناقض كلّه في عبارة «مختزل إلى أبسط صورة»:
فمن دونها لا تناقض عبارةُ «$a$ و $b$ زوجيان» شيئًا.
واِفتراض الاِختزال مشروع لأنّ لكل \omterm{def:g9:fractions:fraction}{كسر} صورةً مختزلة
(\cref{met:g9:arith:simplify}: اِقسم على \omterm{def:g9:arith:gcd}{القاسم المشترك الأكبر})
--- فلو كان $\sqrt2$ كسرًا أصلًا، لكان كسرًا مختزلًا، وذاك
مستحيل.

\textbf{6.} في التفكيك إلى عوامل أولية، تكون المربّعات
زوجية \omterm{def:g8:powers:def}{الأُسُس} (\cref{pb:g9:arith:1}). وفي $a^2 = 3b^2$، يكون
\omterm{def:g8:powers:def}{أُسّ} العدد $3$ زوجيًا في طرف، وفرديًا في الطرف الآخر
(زوجي من $b^2$، مع واحد). ولا عدد له تفكيكان
(\cref{thm:g9:arith:factorization}): فتناقض، والعدد $\sqrt3$ أصمّ.

\textbf{7.} في $a^2 = n b^2$، تجد الحجّة عددًا أوليًا تتناقض
زوجية \omterm{def:g8:powers:def}{أُسّه} بالضبط حين يظهر \omterm{def:g9:arith:prime}{عدد أولي} في $n$
\omterm{def:g8:powers:def}{بأُسّ} \emph{فردي}. وإذا كان كل \omterm{def:g8:powers:def}{أُسّ} في $n$
زوجيًا، كان $n$ مربّعًا تامًّا ($n = m^2$، والعدد
$\sqrt n = m$ صحيح). والنتيجة الكاملة: $\sqrt n$ أصمّ من أجل كل
عدد صحيح $n$ ليس مربّعًا تامًّا ---
$\sqrt2, \sqrt3, \sqrt5, \sqrt6, \sqrt7, \sqrt8, \sqrt{10},
\dots$

\textbf{8.} من $r x = q$ حيث $r \neq 0$: يكون
$x = \frac qr$، وهو \omterm{def:g3:division:remainder}{خارج قسمة} ناطقين، ومن ثمّ ناطق. إذن إذا
كان $x$ أصمّ، فلا يمكن أن يكون $r x$ ناطقًا: فالعددان $2\sqrt2$
و $\frac{\sqrt2}{2}$ أصمّان (والمضروب فيهما $2$
و $\frac12$).

\textbf{9.} لو كان $1 + \sqrt2 = q$ ناطقًا، لكان
$\sqrt2 = q - 1$ ناطقًا: وهذا تناقض. ولو كان
$s = \sqrt2 + \sqrt3$ ناطقًا، لكان
$s^2 = 2 + 2\sqrt6 + 3 = 5 + 2\sqrt6$ ناطقًا أيضًا، فيكون
$\sqrt6 = \frac{s^2 - 5}{2}$ ناطقًا --- لكنّ $\sqrt6$ أصمّ
(السؤال 7). إذن $\sqrt2 + \sqrt3$ أصمّ.

\textbf{10.} المجاميع: العددان $\sqrt2$ و $-\sqrt2$ (أو $\sqrt2$
و $1 - \sqrt2$) كلاهما أصمّ، ومجموعاهما ناطقان $0$ و $1$.
\omterm{def:g2:mult:def}{والجداءات}: $\sqrt2 \times \sqrt2 = 2$. فالصمم قد يتبخّر في
الحساب --- ومن هنا البراهين حالةً حالة.

\textbf{11.} من $x = 1$: متوسّط $1$ و $\frac21 = 2$ هو
$x_1 = \frac32$. ثم متوسّط $\frac32$ و
$\frac{2}{3/2} = \frac43$:
$x_2 = \frac12\left(\frac32 + \frac43\right)
= \frac12 \cdot \frac{17}{6} = \frac{17}{12} \approx 1.4167$.

\textbf{12.} $x_3 = \frac12\left(\frac{17}{12} +
\frac{24}{17}\right) = \frac12 \cdot
\frac{289 + 288}{204} = \frac{577}{408} = 1.4142156\ldots$
وأمام $\sqrt2 = 1.4142135\ldots$: خمسة أرقام عشرية صحيحة بعد ثلاث
خطوات --- فالوصفة تضاعف تقريبًا عدد الأرقام الصحيحة في كل
جولة.

\textbf{13.} إذا كان $x^2 > 2$ (تخمين أكبر من اللازم)، كان
$x > \sqrt2$، و $\frac2x < \frac{2}{\sqrt2} = \sqrt2$: فالضلع
المرافق أصغر من اللازم. وإذا كان $x^2 < 2$، اِنعكست المتباينات.
وفي الحالتين يقع $\sqrt2$ بين $x$ و $\frac2x$، فيكون متوسّطهما ---
وهو التخمين التالي --- أقرب من أسوأ الضلعين: فالمستطيل
ذو \omterm{def:g6:measure:area}{المساحة} $2$ يزداد تربيعًا في كل خطوة.

\textbf{14.} $9 - 8 = 1$؛ \ $289 - 288 = 1$؛ \
$332\,929 - 332\,928 = 1$. فكل تخمين لهيرون $\frac ab$ يحقّق
$a^2 - 2b^2 = 1$: وقد برهن الجزء الأول على اِستحالة بلوغ $0$،
وهذه الكسور تخطئ المستحيل بوحدة واحدة بالضبط --- وهو أقرب ما
تصل إليه الأعداد الصحيحة.

\textbf{15.} يجب أن يحوي \omterm{def:g6:lines:objects}{مستقيم} الأعداد أعدادًا وراء الكسور ---
أطوالًا مثل $\sqrt2$، كتابتها العشرية تمضي أبدًا من غير أن تصير
دورية: وهي \emph{الأعداد الحقيقية}. وبناؤها الأمين حكاية لمجلّد
الثانوية، وبكل صرامة لمجلّدات الجامعة.
\end{solution}
